\section{Phân tích hệ thống và trải nghiệm hiện tại}

\subsection{Hệ thống và quy trình hiện tại}
Hiện nay, hầu hết các cơ sở spa và chăm sóc thú cưng quy mô vừa và nhỏ tại Việt Nam chưa áp dụng hệ thống phần mềm chuyên biệt phục vụ trải nghiệm của khách hàng. Quy trình tương tác giữa chủ nuôi và cơ sở chủ yếu diễn ra qua ba kênh phân mảnh:
\begin{itemize}
    \item \textbf{Kênh gọi điện thoại trực tiếp}: Khách hàng gọi điện tới số hotline của cơ sở để hỏi giờ trống và đặt lịch. Nhân viên tiếp tân phải vừa nghe máy vừa tra cứu sổ tay hoặc lịch bàn thủ công.
    \item \textbf{Kênh nhắn tin qua mạng xã hội (Facebook Fanpage, Zalo, Instagram)}: Khách hàng gửi tin nhắn văn bản hỏi dịch vụ. Do nhân viên tiếp tân thường phải kiêm nhiệm nhiều công việc tại cửa hàng, tin nhắn thường không được phản hồi ngay lập tức.
    \item \textbf{Ghi chép tiếp nhận tại chỗ (Sổ tay giấy, bảng trắng)}: Khi khách mang thú cưng đến tiệm, thông tin về gói dịch vụ và các dặn dò đặc biệt được ghi chép vội vàng lên sổ tiếp nhận hoặc giấy dán ghi chú (\textit{post-it}) dán lên lồng thú cưng.
\end{itemize}

Đặc tính chung của các kênh này là tính \textbf{bất đồng bộ, thiếu liên kết dữ liệu và phụ thuộc hoàn toàn vào trí nhớ cá nhân của nhân viên tiếp nhận}.

\subsection{Phân tích tác vụ}
Dựa trên tài liệu đề xuất đề tài (\textit{Proposal}), quy trình chăm sóc thú cưng hiện tại (Quy trình As-Is) bao gồm 6 bước tác vụ tuần tự:

\begin{enumerate}
    \item \textbf{Bước 1 -- Chủ nuôi gọi điện hoặc nhắn tin hỏi lịch trống}: Khi có nhu cầu, chủ nuôi phải chủ động mở ứng dụng nhắn tin hoặc gọi điện cho tiệm để hỏi xem khung giờ mong muốn (ví dụ: sáng thứ Bảy lúc 09:00) còn nhận thú cưng hay không.
    \item \textbf{Bước 2 -- Chủ nuôi chờ phản hồi và chịu sự thiếu chắc chắn}: Chủ nuôi rơi vào trạng thái chờ đợi tin nhắn phản hồi. Trong nhiều trường hợp, nhân viên trả lời chậm dẫn đến việc khung giờ đó đã kín, hoặc lịch hẹn bị ghi nhận sai lệch do trao đổi qua lại nhiều lần.
    \item \textbf{Bước 3 -- Chủ nuôi khai báo lại tiền sử dị ứng và dặn dò đặc biệt}: Chủ nuôi phải tự mình gõ lại tin nhắn hoặc nói miệng các lưu ý quan trọng (như thú cưng bị dị ứng xà phòng có mùi nồng, da dễ mẩn đỏ, sợ tiếng sấy hoặc nhút nhát với chó dữ).
    \item \textbf{Bước 4 -- Chủ nuôi trực tiếp bàn giao thú cưng tại cơ sở}: Khi đến tiệm, chủ nuôi phải nhắc lại một lần nữa các lưu ý cho nhân viên tiếp tân và hy vọng nhân viên này sẽ truyền đạt đầy đủ cho kỹ thuật viên tắm sấy bên trong.
    \item \textbf{Bước 5 -- Khoảng thời gian chờ đợi mù mờ thông tin}: Trong suốt 2--3 tiếng thú cưng ở tiệm, chủ nuôi hoàn toàn không có bất kỳ cập nhật nào về tiến độ. Nếu muốn biết tình trạng bé, chủ nuôi buộc phải gọi điện chen ngang công việc của tiệm.
    \item \textbf{Bước 6 -- Cơ sở thông báo khi hoàn tất và kết thúc rời rạc}: Cơ sở gọi điện hoặc nhắn tin báo chủ nuôi đến đón. Sau khi thanh toán tiền mặt/chuyển khoản, hóa đơn giấy dễ bị vứt bỏ, thông tin về loại sữa tắm trị liệu đã sử dụng và ghi chú của thợ không được lưu trữ thành lịch sử nhất quán.
\end{enumerate}

\subsection{Luồng người dùng và quy trình tác vụ}
Sơ đồ luồng tương tác hiện tại (Quy trình As-Is) bộc lộ nhiều điểm nghẽn nghiêm trọng (\textit{Bottlenecks}) làm đứt gãy trải nghiệm người dùng (Bảng \ref{tab:flow_bottlenecks}).

\begin{table}[H]
\centering
\small
\begin{tabularx}{\textwidth}{|l|l|X|l|}
\hline
\textbf{Bước tác vụ} & \textbf{Kênh tương tác} & \textbf{Điểm nghẽn trải nghiệm (Bottleneck)} & \textbf{Thời gian trễ} \\ \hline
1. Hỏi lịch & Nhắn tin / Gọi điện & Chờ nhân viên mở sổ đối soát thủ công. & 15--120 phút \\ \hline
2. Xác nhận lịch & Tin nhắn xác nhận & Dễ trùng lịch, ghi sai giờ hoặc lỡ khung giờ đẹp. & Không xác định \\ \hline
3. Khai báo dặn dò & Lời nói / Chat & Phải nhắc đi nhắc lại nhiều lần, thông tin phân mảnh. & 5--10 phút \\ \hline
4. Bàn giao tại tiệm & Trao đổi trực tiếp & Bàn giao qua trung gian tiếp tân, dễ quên khi đổi ca. & 10--15 phút \\ \hline
5. Chăm sóc tại tiệm & Không có kênh cập nhật & Mù mờ thông tin, lo âu, gọi điện làm phiền hai bên. & 120--180 phút \\ \hline
6. Trả bé \& Lưu trữ & Gọi điện đón / Tiền mặt & Thất lạc chi tiết sản phẩm, không có lịch sử tra cứu. & Mất dấu dữ liệu \\ \hline
\end{tabularx}
\caption{Bảng phân tích điểm nghẽn trong luồng tác vụ hiện tại (As-Is)}
\label{tab:flow_bottlenecks}
\end{table}

\subsection{Các vấn đề về trải nghiệm người dùng}
Dưới góc độ khoa học Tương tác Người -- Máy (HCI), các hạn chế của quy trình hiện tại được phân tích thành 4 vấn đề trải nghiệm cốt lõi:

\begin{enumerate}
    \item \textbf{Vấn đề 1 -- Tải trọng nhận thức cao và sự bất định (High Cognitive Load \& Uncertainty)}:
    Việc không thể nhìn thấy trực quan các khung giờ trống khiến người dùng phải liên tục suy đoán, trao đổi đàm phán nhiều vòng qua tin nhắn. Sự chậm trễ trong phản hồi vi phạm nguyên tắc phản hồi tức thì trong tương tác, gây ức chế tâm lý và làm lãng phí thời gian của chủ nuôi bận rộn.
    \item \textbf{Vấn đề 2 -- Nguy cơ sai sót y tế do đứt gãy truyền thông (Communication Breakdown \& Medical Risk)}:
    Dặn dò y tế (dị ứng, bệnh lý, tính cách) là thông tin sinh mệnh đối với thú cưng nhưng lại được truyền đạt qua các kênh phi cấu trúc (lời nói, tin nhắn chat). Khi nhân viên tiếp tân chuyển giao cho kỹ thuật viên tắm sấy, hoặc khi tiệm thay ca trực, thông tin này có xác suất thất lạc lên tới hơn 80\%. Điều này vi phạm nghiêm trọng nguyên tắc an toàn trong thiết kế hệ thống (\textit{Safety \& Error Prevention}) \cite{norman2013design}.
    \item \textbf{Vấn đề 3 -- Vi phạm nguyên lý minh bạch trạng thái hệ thống (Violation of Visibility of System Status)}:
    Theo nguyên lý Heuristic kinh điển của Nielsen \cite{nielsen10heuristics}, hệ thống phải luôn cung cấp phản hồi kịp thời cho người dùng về những gì đang diễn ra. Trong quy trình cũ, khoảng trống thông tin 2--3 tiếng đẩy chủ nuôi vào trạng thái lo lắng tột độ nhưng lại bị ngăn cản liên lạc vì sợ làm phiền. Đây là một điểm mù trải nghiệm (\textit{experience black-box}) nghiêm trọng.
    \item \textbf{Vấn đề 4 -- Thiếu tính cá nhân hóa và khả năng nhận biết thay vì nhớ lại (Violation of Recognition over Recall)}:
    Do không có hệ thống lưu trữ lịch sử có cấu trúc, mỗi lần đến tiệm là một lần người dùng phải bắt đầu lại từ con số không (\textit{cold-start}): tự nhớ lại loại sữa tắm lần trước bé dùng có bị ngứa không, giá tiền bao nhiêu, ngày nào nên đi tắm lại. Việc ép người dùng phải nhớ lại toàn bộ thông tin thay vì hệ thống tự nhận diện và gợi ý vi phạm nguyên lý thiết kế tương tác chuẩn mực \cite{sharp2019interaction}.
\end{enumerate}
